% =============================================================================
%  CHAPTER 5 -- EXPERIMENTAL SETUP                                     [v3 OWNS]
% -----------------------------------------------------------------------------
%  v3.7 (2026-09-14): every section is organised by the three environments, in
%  the order of the study -- obstacle avoidance (the benchmark of DPCC, taken
%  over unchanged), then vision-conditioned alignment and the quadrotor benchmark
%  (both built for this thesis). How each environment is constructed and which
%  constraint forms it uses is v2's Section 4.7 (sec:method:envs); this chapter
%  holds the experimental values, the data, the baselines, the metrics and the
%  protocol. Individual changelog: changelogs/v3.7_20260914_ch5-8_restructure.md
%
%  Vocabulary: Writing/Auxiliary/Naming/TRANSLATION_20260914_dev_jargon_to_scientific.md
% =============================================================================

\chapter{Experimental Setup}\label{ch:setup}

The evaluation follows the two stages of \autoref{sec:intro:problem}. The generative models and
projection methods are first established on the obstacle-avoidance benchmark of \ac{DPCC}, with its
setup taken over unchanged, so that the diffusion baseline is met on its own benchmark. They are then
tested further in two environments built for this thesis: a vision-conditioned alignment task and a
quadrotor benchmark. \autoref{sec:method:envs} describes how each environment is constructed; this
chapter gives, for each of the three, the task and constraint values, the data, the baseline, the
metrics and the evaluation protocol.

% =============================================================================
\section{Tasks and Environments}\label{sec:setup:tasks}
% =============================================================================

\subsection{Obstacle Avoidance}\label{sec:setup:tasks:avoiding}

The obstacle-avoidance task of D3IL \parencite{jia2024towards} is used exactly as \ac{DPCC} uses it
\parencite[Sec.~6.1]{romer2025diffusion} (\autoref{sec:method:env:avoiding}). A manipulator moves its
end effector across a field of six obstacles to a goal line (\autoref{fig:env-avoiding}). At test time
the plan is constrained by halfspaces that the demonstrations never respect, in three geometries --- \emph{top-left-hard}, \emph{top-right-hard} and \emph{both-hard}, defined in \autoref{tab:avoiding-geometries} and drawn in \autoref{fig:constraints-avoiding} --- together with the circular keep-out regions and the action box
of \autoref{sec:method:env:avoiding}. The geometries are those of the \ac{DPCC} configuration. Every
constraint also has a tightened version, shrunk by $0.025$ (\autoref{sec:method:dpcc}).
\srcnote{\texttt{config/projection\_eval.yaml}: \texttt{halfspace\_constraints['avoiding-d3il']},
\texttt{enlarge\_constraints: \{'avoiding': 0.025\}}; parity with \texttt{aux\_repo/dpcc} checked in
v2, \autoref{sec:method:env:avoiding}.}

\begin{table}[htpb]
  \caption[Constraint geometries of obstacle avoidance]{The three constraint geometries of the
  obstacle-avoidance task, in metres. Each halfspace keeps the planned end-effector position below the line
  through the two points; the keep-out disk keeps it outside the circle. All three share the workspace
  $x\in[0.2,0.8]$, $y\in[-0.3,0.4]$, the six obstacles of radius $0.025$ and the goal line $y=0.35$. The
  tightened versions move each line and enlarge each disk by $0.025$. The last column counts the D3IL
  demonstrations whose measured position satisfies the constraints at every step.}
  \label{tab:avoiding-geometries}
  \centering
  \begin{tabularx}{\linewidth}{@{}l L l r@{}}
    \toprule
    Geometry & Halfspace: below the line through & Keep-out disk & Satisfied by \\
    \midrule
    top-left-hard  & $(0.8,\,-0.5)$ and $(0.4,\,0.5)$ & $(0.4,\,0.08)$, $r=0.08$ & 0 of 96 \\
    top-right-hard & $(0.2,\,-0.5)$ and $(0.6,\,0.5)$ & $(0.6,\,0.08)$, $r=0.08$ & 1 of 96 \\
    both-hard      & $(0.8,\,-0.3)$ and $(0.575,\,0.5)$; $(0.2,\,-0.3)$ and $(0.425,\,0.5)$ & $(0.5,\,-0.09)$, $r=0.08$ & 2 of 96 \\
    \bottomrule
  \end{tabularx}
\end{table}
\srcnote{\texttt{config/projection\_eval.yaml} (\texttt{halfspace\_constraints}, \texttt{obstacle\_constraints});
geometry-to-constraint mapping \texttt{scripts/eval.py:93--102}; counts by
\texttt{Data\_Analysis/DA\_in\_Paper/plotting/extract/avoiding\_scene.py}, the check of
\texttt{scripts/visualize\_data\_constraints.py}.}

\begin{figure}[htpb]
  \centering
  \includegraphics[width=0.62\linewidth]{fig_env_avoiding}
  \caption[The obstacle-avoidance task]{The obstacle-avoidance task of D3IL seen from above, with the
  measured end-effector position of its 96 demonstrations. All start at the bottom and cross the goal
  line (green) between the six obstacles (red).}
  \label{fig:env-avoiding}
\end{figure}

\begin{figure}[htpb]
  \centering
  \includegraphics[width=\linewidth]{fig_constraints_avoiding}
  \caption[The three constraint geometries]{The three constraint geometries of
  \autoref{tab:avoiding-geometries}. Shaded: the region excluded by the halfspaces and the keep-out disk;
  dashed: the tightened boundary. The demonstrations (grey) cross the excluded regions; 0, 1 and 2 of the
  96 satisfy the constraints of the three geometries.}
  \label{fig:constraints-avoiding}
\end{figure}

\subsection{Vision-Conditioned Alignment}\label{sec:setup:tasks:aligning}

The alignment task of D3IL \parencite{jia2024towards} asks the manipulator to push a box into a target
pose on a table. Its simulation, demonstrations and vision pipeline are D3IL's. Around them this thesis
built a constrained-control environment (\autoref{sec:method:env:aligning}): a three-dimensional plan
with its dynamics rows, a constraint set on the planned end-effector position and on the action, and
the coupling of the visual encoder to the generative network. The policy observes two RGB camera
streams, an overhead view and an in-hand view, instead of the state (\autoref{fig:env-aligning}). Each
episode starts from an initial box pose, called a context. The tightened constraints are shrunk by
$0.03$\,m.
\srcnote{\texttt{config/visual\_aligning\_eval.yaml}: \texttt{workspace\_bounds},
\texttt{action\_bounds: 'auto'}, \texttt{enlarge\_constraints: 0.03}.}

\hole{Give the workspace box of the planned end-effector position (lower and upper bound per axis) for
each geometry used, read off \texttt{config/visual\_aligning\_eval.yaml}.}

\begin{figure}[htpb]
  \centering
  \todofigure[0.26\textwidth]{The alignment scene with the box and its target pose, next to the two
    camera images the policy receives (overhead and in-hand).}
  \caption[The vision-conditioned alignment task]{The vision-conditioned alignment task built from
  D3IL. The box is pushed into a target pose; the policy observes two RGB streams.}
  \label{fig:env-aligning}
\end{figure}

\subsection{Quadrotor Benchmark}\label{sec:setup:tasks:uav}

The quadrotor benchmark was built for this thesis around the Skydio X2 model of MuJoCo Menagerie
\parencite{zakka2022menagerie} (\autoref{sec:method:env:uav}). The vehicle is underactuated and open-loop
unstable (\autoref{sec:bg:quadrotor}), and its plan is tracked by the controller of
\autoref{sec:method:deployment}. Three scenes are evaluated, in increasing difficulty
(\autoref{fig:env-uav}, \autoref{tab:uav-scenes}):

\begin{description}
  \item[Corridor] --- the entry test. Two parallel walls leave a clear width of $1.90$\,m. A sloping
    halfspace, the slide, descends at $14^\circ$ from the upper wall at the entrance to $0.05$\,m below
    the centre line at the exit, and forces every route to descend by $0.25$--$0.49$\,m before the exit.
    The models are trained on the demonstrations of a narrower corridor and never see the slide, so it is
    a test-time constraint in the same sense as the halfspaces of obstacle avoidance.
  \item[Pillars] --- the main test. Six cylindrical pillars stand in two rows across three columns, and
    the vehicle must choose a way through them.
  \item[S-curve] --- the most difficult test. Two offset corridor segments, joined by a diagonal passage,
    force two sharp changes of direction.
\end{description}

\begin{table}[htpb]
  \caption[Quadrotor scenes]{Geometry of the quadrotor scenes. Walls and pillars enter the projection as
  the constraint forms of \autoref{sec:method:env:uav}, inflated by the vehicle radius of $0.31$\,m.}
  \label{tab:uav-scenes}
  \centering
  \begin{tabularx}{\linewidth}{@{}l L L@{}}
    \toprule
    Scene & Constraints & Values \\
    \midrule
    Corridor & wall halfspaces, $x\in[-2,2]$ & inner faces $y=\pm0.95$\,m (clear $1.90$\,m) \\
             & slide halfspace, $x\in[-2,2]$ & from $(-2.0,\,0.95)$ to $(2.0,\,-0.05)$, $14^\circ$ \\
             & wall ends & disks of radius $0.05$\,m at $x=\pm2$ \\
    Pillars  & six disks & radius $0.12$\,m at $x\in\{-2,0,2\}$, $y=\pm0.6$\,m \\
    S-curve  & switched wall halfspaces & $y\in[-1.25,-0.35]$ for $x\in[-3,-0.5]$; $y\in[0.35,1.25]$ for $x\in[0.5,3]$ \\
             & inner corners & disks of radius $0.05$\,m at $(\mp0.5,\,\mp0.3)$ \\
    All      & workspace box & $x\in[-4,4]$\,m, altitude $[0.30,1.80]$\,m \\
    \bottomrule
  \end{tabularx}
\end{table}
\srcnote{\texttt{config/uav\_projection.yaml}: \texttt{corridor\_v2\_slide} (scene
\texttt{scene\_corridor\_v2.xml}), \texttt{pillars\_hg}, \texttt{s\_curve\_hg};
\texttt{inflation.r\_drone: 0.31}; \texttt{enlarge\_constraints: 0.025}.}

\hole{Confirm that the pillars and s-curve evaluations use the same workspace box as the corridor, and
state the tightening margin ($0.025$\,m in the configuration) for the quadrotor explicitly.}

\begin{figure}[htpb]
  \centering
  \todofigure[0.30\textwidth]{Three panels at a common scale, overhead view, in the order corridor,
    pillars, s-curve. Corridor: both walls, the slide drawn as a line from the entrance to the exit, and
    the three routes L, C, R. Pillars: the six cylinders and the route classes. S-curve: the two
    segments and the diagonal passage. Draw the vehicle to scale (0.62\,m across) in each panel.}
  \caption[The quadrotor scenes]{The three quadrotor scenes, in increasing difficulty: corridor with a
  slide, pillars, and s-curve.}
  \label{fig:env-uav}
\end{figure}

% =============================================================================
\section{Datasets}\label{sec:setup:data}
% =============================================================================

\subsection{Obstacle Avoidance}\label{sec:setup:data:avoiding}

The demonstrations are those of the D3IL obstacle-avoidance task, as used by \ac{DPCC}. Each step
records the commanded and the measured planar end-effector position and the planar increment between
consecutive commanded positions (\autoref{sec:method:env:avoiding}). The dataset holds 96 demonstrations (\autoref{fig:env-avoiding}). None of them satisfies the constraints of top-left-hard, one those of top-right-hard and two those of both-hard (\autoref{tab:avoiding-geometries}), so satisfying the constraints is left to the projection.

\subsection{Vision-Conditioned Alignment}\label{sec:setup:data:aligning}

The demonstrations are those of the D3IL alignment task, with the two camera streams rendered by
D3IL's vision pipeline. For this thesis each step is laid out as the nine-dimensional plan of
\autoref{sec:method:env:aligning} --- the three-dimensional increment and the commanded and measured
end-effector positions --- and paired with the two images.

\subsection{Quadrotor Benchmark}\label{sec:setup:data:uav}

The demonstrations were generated for this thesis as described in \autoref{sec:method:env:uav}. A
reference path through per-scene waypoints, travelled with a cosine speed profile, is flown in MuJoCo by
the cascaded geometric tracking controller, and accepted rollouts are recorded at the planning rate.
\autoref{tab:uav-demos} gives the values per scene. The corridor with the slide reuses the demonstrations
and trained models of the corridor, because the slide is a test-time constraint.

\begin{table}[htpb]
  \caption[Quadrotor demonstrations]{Generation of the quadrotor demonstrations. Shared by all scenes:
  altitude drawn uniformly from $[0.90,1.30]$\,m per episode, physics step $0.01$\,s, recording every
  third physics step ($33$\,Hz), rejection of an episode that falls below $0.50$\,m or touches an obstacle
  on more than the stated fraction of its steps. Durations and offsets are drawn uniformly from the
  stated ranges.}
  \label{tab:uav-demos}
  \centering
  \begin{tabularx}{\linewidth}{@{}l L L L@{}}
    \toprule
    & Corridor & Pillars & S-curve \\
    \midrule
    Passage classes & three lanes at $y\in\{-0.12,\,0,\,0.12\}$\,m & sides of the three pillar pairs: (L,L,L), (L,R,L), (R,L,R), (R,R,R) & one \\
    Reference path & straight, $x$ from $-2.8$ to $2.8$\,m & waypoints joined by arcs of radius $\le0.45$\,m & waypoints joined by arcs of radius $\le0.45$\,m \\
    Duration & $[6,10]$\,s & $[10,16]$\,s & $[16,22]$\,s \\
    Lateral offset & start, $[-0.03,0.03]$\,m, centre lane & --- & whole path, $[-0.04,0.04]$\,m \\
    Maximum contact fraction & $2\,\%$ & $0.1\,\%$ & $8\,\%$ \\
    \bottomrule
  \end{tabularx}
\end{table}
\srcnote{\texttt{uav\_expert\_data\_collect/generator.py} (\texttt{\_build\_traj\_and\_init}: altitude, durations,
offsets; \texttt{SCENE\_MAX\_CONTACT\_FRACTION}; \texttt{Z\_FLOOR\_MARGIN}); \texttt{trajectories.py}
(\texttt{BLEND\_RADIUS}, \texttt{CORRIDOR\_CHANNELS}, \texttt{corridor\_path}); \texttt{dataset\_writer.py}
(\texttt{DATASET\_HZ}); \texttt{quadrotor\_modified.xml} (timestep).}

\hole{Number of accepted demonstrations per scene, and the controller gains of each collection run (the
collector takes the gain set as an argument). Also the number of alignment demonstrations and the
train/validation split of all three environments.}

% =============================================================================
\section{Baselines}\label{sec:setup:baselines}
% =============================================================================

Three generative models are compared with the diffusion model of \ac{DPCC}: flow matching
\parencite{lipman2023flow}, MeanFlow \parencite{geng2025meanflow} and consistency training
\parencite{zhang2025alphaflow}. Two projection methods are compared: the per-step projection of
\ac{DPCC} and endpoint projection \parencite{li2025hardflow}. All generative models share the temporal
U-Net backbone of the baseline, with $4.0$\,M parameters, unless a table states otherwise.

\subsection{Obstacle Avoidance}\label{sec:setup:baselines:avoiding}

The baseline is \ac{DPCC} in the configuration its authors report: the diffusion model at a step budget
of $\nfe=20$, action loss weight $10$, per-step projection on tightened constraints, and selection by
cumulative projection cost over $B=4$ candidate plans \parencite[Sec.~6.1]{romer2025diffusion}.
\autoref{tab:target} places the published result next to its reproduction in this thesis.

\begin{table}[htpb]
  \caption[The DPCC baseline, published and reproduced]{The DPCC baseline on the obstacle-avoidance
  task: the published result and its reproduction in this thesis, on the metrics of
  \autoref{sec:setup:metrics:avoiding}.}
  \label{tab:target}
  \centering
  \begin{tabularx}{\linewidth}{@{}L cc@{}}
    \toprule
    & Published \parencite[Table~1]{romer2025diffusion} & Reproduced \\
    \midrule
    Control steps to the goal           & $69.0 \pm 12.9$ & $69.0$ \\
    Success rate                         & $1.00$          & $0.997$ \\
    Success with constraint satisfaction & $0.98$          & $0.983$ \\
    Control steps with a violation       & $0.0 \pm 0.3$   & $0.03$ \\
    Wall-clock time per control step     & ---             & $0.563$\,s \\
    \bottomrule
  \end{tabularx}
\end{table}

\dataref{temp/2508/batch\_avoiding\_combined\_20260825\_143212; diffusion $\nfe=20$, cumulative
projection cost, tightened; mean per geometry, then across the three geometries}

The reproduction uses more episodes than the published evaluation
(\autoref{sec:setup:protocol:avoiding}), and its \emph{both-hard} geometry was completed for one training
seed rather than five.

\subsection{Vision-Conditioned Alignment}\label{sec:setup:baselines:aligning}

\hole{Baseline configuration for the alignment task.}

\subsection{Quadrotor Benchmark}\label{sec:setup:baselines:uav}

\hole{Baseline configuration for the quadrotor benchmark.}

% =============================================================================
\section{Metrics}\label{sec:setup:metrics}
% =============================================================================

\subsection{Obstacle Avoidance}\label{sec:setup:metrics:avoiding}

The metrics are those of \ac{DPCC} \parencite[Sec.~6.1]{romer2025diffusion}: the number of control
steps to reach the goal, the success rate of reaching the goal, the success rate of reaching the goal
while satisfying the constraints at every step, and the number of control steps on which a constraint
is violated. This thesis adds the wall-clock time per control step.

\subsection{Vision-Conditioned Alignment}\label{sec:setup:metrics:aligning}

Each episode is scored by
\begin{itemize}
  \item the final distance between the box and its target in the table plane, summarised over the
    contexts by its median and its minimum;
  \item the final orientation of the box relative to the target orientation;
  \item whether the box was moved at all;
  \item the number of control steps on which a constraint is violated;
  \item the wall-clock time per control step.
\end{itemize}
When generative models are compared, the distance is measured on plans executed without projection, so
that it reflects the generative model alone; when projection methods are compared, it is measured with
projection.
\srcnote{\texttt{per\_rollout\_detail.csv}: \texttt{context\_final\_xy\_dist},
\texttt{context\_final\_box\_angle\_deg} against \texttt{context\_target\_angle\_deg}, \texttt{frozen},
\texttt{n\_violations}, \texttt{avg\_time\_ms}.}

\subsection{Quadrotor Benchmark}\label{sec:setup:metrics:uav}

Each flight is scored by
\begin{itemize}
  \item whether it is collision-free: no control step on which the vehicle, with its body radius of
    $0.31$\,m, violates a constraint;
  \item two success criteria --- crossing the finish line, and coming within $0.30$\,m of the goal point of
    its route --- each also reported together with collision-free flight as success with constraint
    satisfaction;
  \item the number of control steps to the goal;
  \item the wall-clock time per control step, split into generation and projection.
\end{itemize}

% =============================================================================
\section{Training and Evaluation Protocol}\label{sec:setup:protocol}
% =============================================================================

\subsection{Obstacle Avoidance}\label{sec:setup:protocol:avoiding}

\ac{DPCC} averages its results over five training seeds and ten test episodes for each constraint
geometry \parencite[Sec.~6.1]{romer2025diffusion}. This thesis evaluates the same five training seeds
with twenty episodes per seed, one hundred episodes per configuration and geometry. Results are averaged
per geometry first and then across the three geometries. \autoref{tab:tiers} shows what the larger
number of episodes changes for the baseline.

\begin{table}[htpb]
  \caption[Effect of the number of episodes]{The baseline at ten and at one hundred episodes per
  geometry: same configurations, same code.}
  \label{tab:tiers}
  \centering
  \begin{tabularx}{\linewidth}{@{}L cc@{}}
    \toprule
    Success with constraint satisfaction & 10 episodes & 100 episodes \\
    \midrule
    Cumulative projection cost, tightened, top-right-hard & 1.00 & 0.95 \\
    Temporal consistency, tightened, top-right-hard       & 1.00 & 0.92 \\
    Untightened configurations                            & \multicolumn{2}{c}{change by up to 0.35} \\
    Tightened configurations                              & \multicolumn{2}{c}{change by at most 0.08} \\
    \bottomrule
  \end{tabularx}
\end{table}

\dataref{DA\_20260819 --- baseline at 10 against 100 episodes per geometry}

At ten episodes a success rate of $1.00$ resolves to $\pm0.10$, and the baseline's perfect rates at ten
episodes become $0.95$ and $0.92$ at one hundred.

\provisional{The flow-based models have not yet been evaluated at the published number of episodes on
the current evaluation code. Until that run is done, the low-budget baseline results in
\autoref{sec:res:fewstep} are at ten episodes and the flow-based models at one hundred.}

The flow-matching results were produced by the sampler inherited from the baseline, which draws the
initial noise at half the unit scale while flow matching was trained at unit scale
(\autoref{sec:method:engine}). MeanFlow and consistency training use unit scale at both ends.
\provisional{The flow-matching models have not been re-evaluated at unit scale.}

\subsection{Vision-Conditioned Alignment}\label{sec:setup:protocol:aligning}

The evaluation resets the environment to a fixed, enumerated set of contexts, so every model and
projection method is evaluated on the same initial box poses and comparisons are paired. The results
use training seed 6 and ten contexts with one episode each, and differences are tested with exact
paired sign tests and paired sign-flip permutation tests. At ten pairs the smallest two-sided sign-test
$p$-value is $0.002$.
\srcnote{\texttt{config/visual\_aligning\_eval.yaml:44--45} (\texttt{n\_contexts},
\texttt{n\_trajectories\_per\_context}); \texttt{eval\_mix\_visual\_aligning.py:1240}.}
\provisional{D3IL evaluates alignment on sixty contexts; the evaluation on fifty contexts is pending.}

\subsection{Quadrotor Benchmark}\label{sec:setup:protocol:uav}

All quadrotor results use training seed 6. The corridor is flown twelve times per configuration, four
times on each of the routes L, C and R; pillars and s-curve are flown ten times per configuration. The
scenes are evaluated and reported separately. The corridor evaluation constrains the commanded position
in addition to the measured one and releases the action box, so its projection configuration differs
from the one used on pillars and s-curve, and the results are not pooled across scenes. The plan is
tracked by the cascaded geometric tracking controller with brake-to-rest velocity setpoints
(\autoref{sec:method:deployment}). The diffusion model on the quadrotor is trained with an action loss
weight of $1$.
\srcnote{\texttt{Gen15/U16/DA\_20260914\_corridor\_v2\_paper\_full.md} \S1 (\texttt{-bounds\_free-pdes-tightened},
\texttt{pid\_stopgo}, $n=12$); \texttt{Campaign\_20260907\_five\_missions/DA\_20260912\_pillars\_diffusion\_baseline\_reference.md} \S8.}

\hole{Confirm the controller and velocity-setpoint policy used for the pillars and s-curve results, and
add seeds once the multi-seed quadrotor evaluation is run.}

\subsection{Compute Environment}\label{sec:setup:protocol:compute}

All training and evaluation reported in this thesis ran on a single shared GPU node of the
institute's cluster, scheduled with SLURM. The node carries two 16-core AMD EPYC~7282 processors
(32~physical cores, 64~hardware threads) and eight NVIDIA RTX~A5000 accelerators with 24~GB of
device memory each. \textbf{Every job was allocated exactly one GPU}: no result in this thesis comes
from a distributed or multi-GPU run. The full hardware and software inventory, including package
versions, is given in \autoref{app:repro}, \autoref{tab:compute}.
\srcnote{Slurm\_Codes/sbatch/templates/2026\_04\_30\_job\_template.sh; every log header in Slurm\_Codes/logs/}

Two properties of this machine bear on how the wall-clock measurements are read. The node is shared:
it serves several GPU queues at once, but the scheduler neither oversubscribes GPUs nor preempts
running jobs, so a job holds its accelerator exclusively and contention is confined to the CPU, the
memory system and I/O. And thirty-two physical cores serve eight accelerators, while both the MuJoCo
rollout and the per-step projection, solved by SLSQP (\autoref{sec:method:proj}), run on the CPU, so
evaluation throughput is plausibly limited by the CPU rather than the RTX~A5000. Wall-clock times are
therefore compared between configurations measured on the same node, and not read as hardware
benchmarks.
